\documentclass[10pt,a4paper,ragged2e,withhyper,paragraphstrue]{altacv}

% \newenvironment{sloppypar*}{\sloppy\ignorespaces}{\par}
% \geometry{left=1.2cm,right=1.2cm,top=1cm,bottom=1cm,columnsep=0.75cm}
% \geometry{left=2.54cm,right=2.54cm,top=2.54cm,bottom=2.54cm,columnsep=0.75cm}
\geometry{left=2.0cm,right=2.0cm,top=2.0cm,bottom=2.0cm,columnsep=0.75cm}
\usepackage{subfiles} % Use subfiles to define personal information.
\usepackage{paracol}

% Add technical terms to dictionary
\usepackage{xspace}

% % Change the font if you want to, depending on whether
% % you're using pdflatex or xelatex/lualatex
% \ifxetexorluatex
%   % If using xelatex or lualatex:
%   \setmainfont{Roboto Slab}
%   \setsansfont{Lato}
%   \renewcommand{\familydefault}{\sfdefault}
% \else
%   % If using pdflatex:
%   \usepackage[rm]{roboto}
%   \usepackage[defaultsans]{lato}
%   % \usepackage{sourcesanspro}
%   \renewcommand{\familydefault}{\sfdefault}
% \fi

% Fork (before v1.6.5a): Change the color codes to test your personal variant on any mode
\ifdarkmode%
  \definecolor{PrimaryColor}{HTML}{C69749}
  \definecolor{SecondaryColor}{HTML}{D49B54}
  \definecolor{ThirdColor}{HTML}{1877E8}
  \definecolor{BodyColor}{HTML}{ABABAB}
  \definecolor{EmphasisColor}{HTML}{ABABAB}
  \definecolor{BackgroundColor}{HTML}{191919}
\else%
%  \definecolor{PrimaryColor}{HTML}{001F5A}
%  \definecolor{SecondaryColor}{HTML}{0039AC}
%  \definecolor{ThirdColor}{HTML}{F3890B}
%  \definecolor{BodyColor}{HTML}{666666}
%  \definecolor{EmphasisColor}{HTML}{2E2E2E}
%  \definecolor{BackgroundColor}{HTML}{E2E2E2}
  \definecolor{PrimaryColor}{HTML}{001F5A}
  \definecolor{SecondaryColor}{HTML}{448888}
  \definecolor{ThirdColor}{HTML}{4682B4}
  \definecolor{BodyColor}{HTML}{666666}
  \definecolor{EmphasisColor}{HTML}{2E2E2E}
  \definecolor{BackgroundColor}{HTML}{FFFFFF}
%  \definecolor{BackgroundColor}{HTML}{E2E2E2}
\fi%

%My Colours:
% Teal 448888
% Orange DD5522 
% Light blu 77EEEE
% Steel Blu 4682B4

\colorlet{name}{PrimaryColor}
\colorlet{tagline}{SecondaryColor}
\colorlet{heading}{PrimaryColor}
\colorlet{headingrule}{ThirdColor}
\colorlet{subheading}{SecondaryColor}
\colorlet{accent}{SecondaryColor}
\colorlet{emphasis}{EmphasisColor}
\colorlet{body}{BodyColor}
\pagecolor{BackgroundColor}

%\usepackage{hyperref
%%Setup bib hyperlinks
\hypersetup{%
	urlcolor=SecondaryColor,
	citecolor=SecondaryColor
}

% Change some fonts, if necessary
\renewcommand{\namefont}{\Huge\rmfamily\bfseries}
\renewcommand{\personalinfofont}{\small\bfseries}
\renewcommand{\cvsectionfont}{\LARGE\rmfamily\bfseries}
\renewcommand{\cvsubsectionfont}{\large\bfseries}

% Change the bullets for itemize and rating marker
% for \cvskill if you want to
\renewcommand{\itemmarker}{{\small\textbullet}}
\renewcommand{\ratingmarker}{\faCircle}

\newcommand{\textalignment}{
	\justifying
    \tolerance=1 %https://tex.stackexchange.com/a/5039/137109
    \emergencystretch=\maxdimen
    \hyphenpenalty=10000 
    \hbadness=10000
}

\setlength{\parindent}{24pt}
\newcommand{\pind}{\hspace{24pt}}

\begin{document}
    
    \IfFileExists{./cl-intro.info}{
            \subfile{./ms-intro.info}
        }{
            \subfile{./ms-intro-example.info}
        }

    % \vspace{1em}
    {\textalignment
    
    % https://www.academictransfer.com/en/blog/writing-a-research-vision-document/
    % Mission statement blog: 

    % Q1. Academic Identity: Before writing your research vision, start by thinking about yourself as a researcher. What are you passionate about? What is the change you wish to see in the world? What do you enjoy doing? Which methods and skills are you good at
    % A1. I am passionate about computation aspects of experimental science, including synthesis and characterisation of electronic materials. I hope to see computational abilities developed that can be utilised in industry to create new types of devices.

    %  I use algorithms, automation and computation to study electronic materials, such as 2D crystal materials, topological materials, and organic semiconductors, in order to demonstrate demonstrate how new materials can be utilised to build new devices and technologies.

    My research is driven by a fundamental intrigue in novel electronic materials, with the goal of demonstrating reliable processability and physical properties for modern electronic devices. Achieving this requires optimising complex synthesis and control systems for in-situ characterisation. This interest stems from my background in condensed matter physics, electronic materials-science engineering, and scientific software development, where I have repeatedly encountered the same core challenge: modern experiments generate vast, complex datasets and rely on intricate instrumentation; however, much of the device lifecycle remains manual—sample fabrication, controlled measurement, analysis, optimisation, and iteration. My central goal is to bridge this gap through automation and, ultimately, autonomous scientific discovery. In particular, I aim to control the creation of complex material systems with novel electronic properties and autonomously explore their physical phase space.

    \pind During my PhD and postdoctoral work in electronic materials, I developed strong expertise across synchrotron, neutron, and nanofabrication facilities. My research includes crystal growth and transfer methods, scattering characterisation, and cryogenic electronic transport, all relying on diverse control systems to study layered 2D systems, topological materials, and organic semiconductors. These systems share a common theme: they require complex, multi-dimensional synthesis and control to achieve desired properties. For example, layered 2D crystals exhibit strong surface sensitivity, often precluding vapour deposition due to damage, and requiring clean transfer methods for state-of-the-art electronic structures. In organic semiconducting polymers, structure-property relationships emerge from a highly coupled processing space, where blade coating, spin coating, and thermal annealing define a high-dimensional parameter space governing morphology, crystallinity, and device performance.

    \pind These materials exemplify key challenges in modern experimental science: relationships between processing conditions and resulting properties are non-linear, history-dependent, and difficult to probe directly, motivating automated approaches. My work using synchrotron X-ray scattering integrates in-situ measurements during solution processing to capture structural evolution in real time, generating rich, multi-dimensional datasets that encode this phase space. However, extracting insight and guiding experimental decisions remains largely manual and intuition-driven, reinforcing my interest in automation as essential for navigating complex experimental landscapes.

    \pind A key component of my postdoctoral work has been developing modern software pipelines and scientific tools (e.g., waveguiding simulations for X-ray scattering, and NEXAFS and Kramers-Kronig analysis toolkits) that are robust, reusable, and aligned with modern software engineering practices. While not control systems themselves, these tools are a step toward more trustworthy, integrated, and automated workflows that can be readily adopted and extended. Such modular approaches (particularly in shared facilities) are essential for consistency and accelerating progress. These efforts reduce manual intervention and enable more systematic exploration of parameter spaces across varying equipment, allowing researchers to focus on higher-level insight.

    \pind Beyond automation, I am motivated by autonomous discovery: systems that iteratively explore high-dimensional experimental spaces and optimise measurement strategies in time-intensive environments, identifying phases and optimal conditions with minimal human input. The processing landscape of polymer materials provides a compelling testbed, where large phase spaces are accessible but difficult to navigate efficiently. These ideas extend naturally to layered 2D materials involving growth and transfer, and more broadly across experimental science and engineering.

    \pind \pind Looking forward, I aim to develop integrated control and analysis systems that unify material synthesis, in-situ characterisation, and data-driven decision-making within a single experimental loop. This includes automated laboratories and adaptive measurement platforms capable of interacting with complex physical systems. More broadly, I am interested in how autonomous systems can navigate complex phase spaces; across materials science, manufacturing, and other domains where high-dimensional optimisation is central. Ultimately, my vision is to transition scientific practice from manual workflows to integrated, autonomous systems that accelerate discovery, where scientific systems become active participants in the discovery process.
    }

    {\color{emphasis}
    \textalignment
    % \vspace{1em}
    % Yours Sincerely, \newline
    % {\color{emphasis}Matthew G Gebert}

    }
    
    % \vspace*{0.5em}
    \divider

    % Use coverletter personal info if it exists, otherwise use example personal info, as the material should be the same as the 
    \IfFileExists{./cl-personal.info}{
        \vspace{-2em}
        \subfile{./cl-personal.info}
    }{
        \vspace{-2em}
        \subfile{./cl-personal-example.info}
    }

\end{document}